% se puede agregar la opción [english] para 
%  memorias o tesis en inglés (borrando el archivo .aux)
\documentclass{umemoria} 

\depto{Departamento de Ciencias de la Computación}
\author{Marcelo Esteban Becerra Alarcón}
\title{Mejora de un sistema de mesa de ayuda para estudiantes del DCC}

% incluir ambos comandos para una doble titulación
%  o quitar el comando que no aplica
\memoria{Ingeniero Civil en Computación}
% \tesis{Magíster en ???}
%\tesis{Doctor en ???} % incluir solo este comando para doctorados

% puede haber varios profesores guía seperados por coma;
% pero si es una memoria, solo puede haber un profesor guía
\guia{Jocelyn Simmonds} 

% puede haber varios profesores co-guía seperados por coma;
% pero si es una memoria, el profesor co-guía será el primer
% integrante de la comisión
%\coguia{Nombre Completo Co-Guía} % incluir en caso de co-guía de *tesis*

%\cotutela{Nombre Institución} % incluir en caso de cotutela

\comision{}

%\auspicio{Nombre Institución} % incluir en caso de recibir financiamiento

% tiene que ser el año en que se da el examen de título/grado (defensa)
%\anho{2021} % incluir solo para reemplazar el año actual

\usepackage{lipsum}

\begin{document}

\frontmatter
\maketitle

\begin{resumen}
   \par El \acrfull{dcc}, ha tenido un aumento sustancial en la cantidad de alumnos que ingresan a la carrera. Esto implica que se generen una serie de desafíos en torno a los procesos académicos. Estos desafíos generan un aumento en la carga de los funcionarios y docentes, a su vez que hacen que no se pueda responder con la velocidad necesaria que requieren los cambios. Esto aumenta la incertidumbre por parte de los alumnos, sobre todo en la forma de abordar los diferentes desafíos de los procesos académicos. Este proyecto busca continuar la extensión de un sistema de mesa de ayuda, pensado para modernizar el trato efectivo entre alumnos, profesores y funcionarios. Mejorando su infraestructura y añadiendo funcionalidades clave para la continuidad y mejora del proyecto.

   \par Este proyecto se realizó en 4 fases. En primer lugar, se efectuó un nuevo levantamiento de datos para evaluar la percepción de los alumnos del sistema actual, y adquirir nociones que permitieran hacer un diseño centrado en el usuario de las nuevas funcionalidades. A partir de estos resultados se integraron los objetivos y preferencias de los alumnos de manera explícita en el sistema, y se agruparon en categorías funcionales y cualitativas las valoraciones de los alumnos.

   \par Luego se hizo un proceso de diseño de alto nivel a través del lenguaje de modelación \gls{i*}. Este se traspasó a un diseño basado en casos de uso a través de \acrshort{uml}. Finalmente, se produjo un diseño de funcionalidades de 3 partes que contempla: las preferencias y objetivos de los usuarios, las opciones en el sistema y, las funcionalidades específicas que permiten dar cumplimento a las dichas alternativas.

   \par Después, se procedió a un proceso de análisis y reestructuración del código actual, lo que permitió añadir nuevas funcionalidades, así como mejorar el sistema existente, asegurando su extensibilidad. Se solucionaron problemas de compatibilidad y \textit{bugs} en el código.

   \par Finalmente, se implementaron nuevas tecnologías como \textit{\gls{Celery}}, que permitieron la implementación de funcionalidades de suscripción personalizada. Permitiendo a cada alumno agregar recordatorios de los procesos habilitados.

   \par Este trabajo extiende las funcionalidades anteriores y reestructura el código existente, de manera de crear un sistema extensible y personalizable, favoreciendo la continuidad de este servicio. Logrando así los objetivos planificados.
   
\end{resumen}

% opcional: incluir para tesis en inglés;
%  en este caso hay que tener el resumen y abstract
%   en ambos idiomas
%\begin{abstract}
%\lipsum[1-4]
%\end{abstract}

\begin{dedicatoria}
Una dedicatoria corta.
\end{dedicatoria}

\begin{thanks}
\lipsum[1-2]
\end{thanks}

\tableofcontents
\listoftables % opcional
\listoffigures % opcional

\mainmatter

% \chapter{Introducción}

\lipsum[30-35]
\begin{enumerate}
	\item Item 1
	\begin{enumerate}
		\item Subitem 1
		\item Subitem 2 (ver Figura \ref{logofcfm})
	\end{enumerate}
	\item Item 2
	\item Item 3
\end{enumerate}
\begin{teo}
Se tiene que $$\int_0^t e^sds=e^t-1.$$
\end{teo}
\lipsum[36-40]
% \chapter{Segundo}

\lipsum[1-3]

\begin{defn}[ver \cite{KAR00}] Definición definitiva $$\frac{d}{dx}\int_a^xf(y)dy=f(x).$$\end{defn}
% \chapter{Tercero}
\lipsum[50-60]
% \chapter{Conclusión}

\lipsum[130-132]
\begin{figure}
	\centering
	\includegraphics[scale=.2]{imagenes/fcfm.pdf}
	\caption{Logo de la Facultad}
	\label{logofcfm}
\end{figure}

\lipsum[133-134]

\begin{table}
	\centering
	\caption{Tabla 1}
	\label{tabla:1}
	\begin{tabular}{|c|c|r|}
		\hline
		\textbf{Campo 1} & \textbf{Campo 2} & \textbf{Num} \\\hline
		Valor 1a & Valor 2a & 3\\
		Valor 1b & Valor 2b & 3\\
		\hline
	\end{tabular}

\end{table}

\lipsum[135]


% ver https://www.overleaf.com/learn/latex/Glossaries
% 
\newglossaryentry{latex}
{
    name=latex,
    description={Is a mark up language specially suited for scientific documents}
}

\newacronym{cadcc}{CADCC}{Centro de Alumnos del Departamento de Ciencias de la Computación}

\newacronym{dcc}{DCC}{Departamento de Ciencias de la Computación, de la Universidad de Chile}

\newacronym{ps}{CC5402}{CC5402 Proyecto de software} 

\newacronym{e}{CC6908}{CC6908 Introducción al trabajo de título}

\newacronym{f1}{CC6909}{CC6909 Trabajo de Titulo}

\newacronym{f2}{CC6910}{CC6910 Trabajo de Memoria de Titulo}

\newacronym{S}{\textit{Alumno S}}{alumno auto suficiente}

\newacronym{R}{\textit{Alumno R}}{alumno regular}

\newacronym{P}{\textit{Alumno P}}{alumno práctico}

\newacronym{UCD}{UCD}{User Centered Desing}

\newacronym{f}{FG}{\textit{Focus Group}}

\newglossaryentry{i*}{
    name=i*,
    description={El lenguaje de modelamiento i*, fue introducido como un marco de referencia orientado a modelar actores y objetivos. Consiste en un lenguaje de modelamiento junto con una serie de tecnicas que permiten analizar estos modelos. \cite{Dalpiaz2016}}
    }
    
\newglossaryentry{effectiveGUI}{
    name=Efective GUI,
    description= {Cuando se habla de \textit{Efective GUI} o interfaz gráfica efectiva para el usuario, se habla de una interfaz que está diseñada para simular o reemplazar una interacción humana, se acuña principalmente en el área de IA relacionada con los chatbots}
    }

\newglossaryentry{Telegram}{
    name=Telegram,
    description= {Telegram es un software gratiuto, multiplataforma, basado en la nube, de mensajería instantanea. El servicio además provee video encriptado de principio a fin, VoIP, envío de archivos y varias otras funciones. Fue lanzado para iOS en Agosto 14 2013 y Android en Octubre 2013 }
}

\newglossaryentry{Celery}{
    name=Celery,
    description= {Telegram es un software gratiuto, multiplataforma, basado en la nube, de mensajería instantanea. El servicio además provee video encriptado de principio a fin, VoIP, envío de archivos y varias otras funciones. Fue lanzado para iOS en Agosto 14 2013 y Android en Octubre 2013 }
}

\newglossaryentry{Django}{
    name=Django,
    description= {Telegram es un software gratiuto, multiplataforma, basado en la nube, de mensajería instantanea. El servicio además provee video encriptado de principio a fin, VoIP, envío de archivos y varias otras funciones. Fue lanzado para iOS en Agosto 14 2013 y Android en Octubre 2013 }
}

\newglossaryentry{ORM}{
    name=ORM,
    description= {Telegram es un software gratiuto, multiplataforma, basado en la nube, de mensajería instantanea. El servicio además provee video encriptado de principio a fin, VoIP, envío de archivos y varias otras funciones. Fue lanzado para iOS en Agosto 14 2013 y Android en Octubre 2013 }
}

\newglossaryentry{ODM}{
    name=ODM,
    description= {Telegram es un software gratiuto, multiplataforma, basado en la nube, de mensajería instantanea. El servicio además provee video encriptado de principio a fin, VoIP, envío de archivos y varias otras funciones. Fue lanzado para iOS en Agosto 14 2013 y Android en Octubre 2013 }
}

\newglossaryentry{REST}{
    name=REST,
    description={Representational State Transfer, architectural style for distributed hypermedia systems}
}
 % opcional

\nocite{*}
\bibliographystyle{plain}
\bibliography{bibliografia}

% opcional ...
% \begin{appendices}
% % \chapter{Anexo}
\lipsum[50-60]
%  \par sometthing in the appendiex
% \end{appendices}

\end{document}
